\documentclass[11pt,letterpaper]{article}
\usepackage[T1]{fontenc}
\usepackage[utf8]{inputenc}
\usepackage{lmodern}
\usepackage[margin=1in]{geometry}
\usepackage{amsmath,amssymb,amsthm,mathtools}
\usepackage{mathrsfs}
\usepackage{microtype,booktabs,array,longtable,enumitem}
\usepackage[dvipsnames]{xcolor}
\usepackage{hyperref}
\usepackage{fancyhdr}
\hypersetup{colorlinks=true,linkcolor=MidnightBlue,citecolor=MidnightBlue,urlcolor=MidnightBlue,
 pdftitle={Compact-Curve Geometry over Surcomplex Hahn Fields},
 pdfauthor={Research manuscript prepared with ChatGPT},
 pdfsubject={Picard classification, valuation monodromy, Hahn cohomology, and Abel criteria}}
\setlength{\headheight}{14pt}
\pagestyle{fancy}\fancyhf{}
\fancyhead[L]{\small Compact-curve geometry over surcomplex Hahn fields}
\fancyhead[R]{\small Research manuscript}
\fancyfoot[C]{\thepage}
\setlength{\parskip}{3pt}
\setlength{\parindent}{1.1em}
\emergencystretch=2em
\numberwithin{equation}{section}
\newtheorem{theorem}{Theorem}[section]
\newtheorem{proposition}[theorem]{Proposition}
\newtheorem{lemma}[theorem]{Lemma}
\newtheorem{corollary}[theorem]{Corollary}
\theoremstyle{definition}
\newtheorem{definition}[theorem]{Definition}
\newtheorem{example}[theorem]{Example}
\theoremstyle{remark}
\newtheorem{remark}[theorem]{Remark}
\newcommand{\C}{\mathbb C}
\newcommand{\R}{\mathbb R}
\newcommand{\Z}{\mathbb Z}
\newcommand{\Q}{\mathbb Q}
\newcommand{\No}{\mathbf{No}}
\newcommand{\HH}{\mathcal H}
\newcommand{\RR}{\mathcal R}
\newcommand{\PP}{\mathcal P}
\newcommand{\MM}{\mathscr M}
\newcommand{\OO}{\mathcal O}
\newcommand{\Acal}{\mathcal A}
\newcommand{\mm}{\mathfrak m}
\newcommand{\supp}{\operatorname{supp}}
\newcommand{\Pic}{\operatorname{Pic}}
\newcommand{\Piclf}{\operatorname{Pic}_{\mathrm{lf}}}
\newcommand{\Div}{\operatorname{Div}}
\newcommand{\Prin}{\operatorname{Prin}}
\newcommand{\Res}{\operatorname{Res}}
\newcommand{\rank}{\operatorname{rank}}
\newcommand{\im}{\operatorname{im}}
\newcommand{\coker}{\operatorname{coker}}
\newcommand{\Hom}{\operatorname{Hom}}
\newcommand{\Exp}{\operatorname{Exp}_{H}}
\newcommand{\Log}{\operatorname{Log}_{H}}
\newcommand{\db}{\bar\partial}
\newcommand{\Hahn}[1]{#1((t^\Gamma))}
\newcommand{\Hp}{\HH_{\Gamma,+}}
\newcommand{\Hg}{\HH_\Gamma}
\newcommand{\Rg}{\RR_\Gamma}
\newcommand{\Mg}{\MM_\Gamma}
\newcommand{\Kg}{K_\Gamma}
\newcommand{\rg}{R_\Gamma}
\newcommand{\mg}{\mm_\Gamma}
\newcommand{\vclass}{\nu}
\newcommand{\cl}{\operatorname{cl}}
\newcommand{\id}{\operatorname{id}}
\newcommand{\ot}{\otimes_{\C}}
\newcommand{\eps}{\varepsilon}
\newcommand{\tinyurl}[1]{{\small\url{#1}}}

\title{\textbf{Compact-Curve Geometry\\over Surcomplex Hahn Fields}\\[0.65em]
\large Picard classification, valuation monodromy,\\
finite cohomology models, and an exact Abel criterion}
\author{A research manuscript prepared for Vladimir Reshetnikov\\
\small with ChatGPT}
\date{21 September 2026}

\begin{document}
\maketitle
\thispagestyle{empty}
\begin{abstract}
Fix a compact connected ordinary Riemann surface $X$ of genus $g$ and a
set-sized ordered subgroup $\Gamma$ of the surreal numbers. We study the
ordinary-base sheaf of holomorphic-coefficient Hahn series
$\Hg$, with scalar field $\Kg=\C((t^\Gamma))\subset\No[i]$.
This is not the fine topology of the surcomplex plane. We prove an exact
classification of locally free rank-one modules:
\[
 \Piclf(X,\Hg)\simeq H^1(X,\underline\Gamma)\oplus\Pic(X)
       \oplus\bigl(H^1(X,\OO_X)\ot\mg\bigr).
\]
The first summand is precisely the obstruction to the existence of a
nonzero coherent meromorphic section. In its kernel, an explicit finite
matrix, obtained by a support-controlled elimination of the ordinary
Dolbeault complex, computes both cohomology groups and proves
Riemann--Roch without any rank-one valuation hypothesis.
For finitely many infinitesimally moving divisor clusters we give a
necessary and sufficient principality criterion: the ordinary divisor
must be principal and a root-free Hahn-valued Abel functional must vanish.
A genus-two example exhibits a nonlinear obstruction after the first
symmetric displacement cancels.
Finally, for the integral Hahn sheaf we compute exactly what all powers
of one valuation scale can see. The missing Picard classes are
$H^1(X,\OO_X)\ot\bigcap_{n\ge1}t^{n\delta}\rg$; they can be nonzero
at higher surreal ranks. Every infinite construction carries an explicit
well-ordered support certificate.
\end{abstract}
\vspace{0.5em}
\noindent\textbf{Status and scope.}
The manuscript supplies proofs of the statements in the category defined
below. It proposes the combined support-controlled classification,
obstruction, and comparison results as research contributions; it does
not claim to settle a named published conjecture or to certify publication
priority. Classical Hahn algebra, Hodge theory, Picard theory,
Riemann--Roch, Abel's theorem, and homological perturbation are credited
as inputs, not claimed as new. The proofs have not been independently
refereed or proof-assistant checked. Exact finite symbolic checks are
included separately.
\clearpage
\tableofcontents
\clearpage

\section{The question, the category, and the contributions}
\label{sec:intro}

Replacing complex scalars by surcomplex scalars is not a single operation.
One can change the topology, enlarge an algebraic coefficient field,
introduce a valuation, or require well-ordered families of ordinary
holomorphic functions. Those choices lead to different theories. Here we
keep an ordinary compact complex curve as the base space and enlarge its
coefficient sheaf by Hahn series. This choice makes it possible to ask a
precise global question:
\begin{quote}
Which locally trivial one-dimensional Hahn modules are described by
meromorphic divisor data, how are their sections obstructed, and which
obstructions survive every truncation at one infinitesimal scale?
\end{quote}
The answer separates three pieces of data that should not be confused:
valuation monodromy, an ordinary complex line bundle, and an infinitesimal
cohomology class.

\subsection{Relation to the supplied repository}
The repository was inspected at commit
\texttt{aa846271b4dcae2c055b216126a87210292ec19b}.
Its documentation includes common-domain Hahn analytic geometry,
finite deformations, global divisor problems over the ordinary complex
plane, and separate work on differential equations, rank-one analytic
spaces, and finite-dimensional spectral theory~\cite{repo}.
The global-divisors report explicitly emphasizes common ordinary domains
and noncompact support obstructions. It does not provide the compact-base
classification or the fixed-scale Picard comparison proved here.

The present development does not assume the repository's unrefereed
claims about principal ideal stalks, nor identify locally free modules
with a potentially larger category of invertible objects on an arbitrary
ringed space. We work with $\Piclf$ throughout. Every ring, topology, and
notion of meromorphic section used in a theorem is specified below.

\subsection{The main results at a glance}
Let $X$ have genus $g$, let $\Kg=\C((t^\Gamma))$, and let
$\mg$ be its positive-valuation ideal. The main statements are as follows.

\paragraph{Classification and meromorphic obstruction.}
Theorem~\ref{thm:pic} gives a canonical splitting into
$H^1(X,\underline\Gamma)$, $\Pic(X)$, and
$H^1(X,\OO_X)\ot\mg$. Theorem~\ref{thm:meromorphic} shows that a
nonzero coherent meromorphic section exists exactly when the first
component vanishes. Thus the category has high-degree line bundles
without any nonzero holomorphic section. This is a boundary on the
category, not a counterexample to ordinary algebraic Riemann--Roch.

\paragraph{A finite and exact cohomology model.}
For valuation class zero, let $L_0$ be the ordinary component and let
$\alpha$ be a positive-support harmonic $(0,1)$-form encoding the
infinitesimal component. With ordinary Hodge operators $p,i,h$, the matrix
\[
 B_\alpha=p_1\alpha(1+h\alpha)^{-1}i_0
\]
has entries in $\mg$ and computes both cohomology groups
(Theorem~\ref{thm:finite}). Its inverse operator is a Hahn identity,
not a convergent sequence. This gives Riemann--Roch, a sharp large-degree
section count, and an integral model that also records torsion.

\paragraph{An exact Abel test without choosing roots.}
For finite moving clusters, Theorem~\ref{thm:abel} gives an if-and-only-if
criterion using ordinary principality and the power sums of the local
cluster polynomials. The formula is valid in arbitrary ordered groups;
its support proof uses the symmetric coefficients rather than separately
chosen roots. A completely explicit genus-two obstruction appears in
Section~\ref{sec:examples}.

\paragraph{What one-scale completion forgets.}
For $\delta>0$, let $\Delta$ be the convex subgroup generated by $\delta$.
Theorem~\ref{thm:adic} computes the inverse limit of the integral
Picard groups modulo $t^{n\delta}$ and the kernel of the comparison map.
If $\Delta\ne\Gamma$, all powers of that one scale miss genuine classes.
For $\Gamma=\Q+\Q\omega$, a deformation beginning at $t^\omega$ is a
concrete witness.

\subsection{Novelty boundaries and nearby literature}
The relevant classical foundations are generalized power-series algebra
and surreal normal forms~\cite{neumann,gonshor,mantova}, and the analytic
theory of compact Riemann surfaces~\cite{narasimhan}. The finite-matrix
construction is a two-term instance of homological perturbation, not a
new discovery of that lemma~\cite{crainic}. Higher cohomological
obstructions for ordinary line-bundle deformations have a substantial
history, including Green--Lazarsfeld~\cite{greenlazarsfeld}.

Hahn-meromorphic analysis in the sense of M\"uller--Strohmaier involves
convergent expansions near zero on a logarithmic cover~\cite{muller};
our formal coefficient sheaf is different. Logarithmic Picard theory and
its tropicalization already study monodromy and divisorial
representability in a different geometric setting~\cite{molchowise}.
In particular, the existence of a monodromy obstruction is not asserted
to be a new general idea. Here the base is a fixed smooth complex curve,
the relevant group is $H^1$ of its ordinary topology, and no logarithmic
structure, tropical dual graph, or bounded-monodromy condition is imposed.

The proposed contribution is the precise combination of these questions
in the common-domain Hahn category, including arbitrary-rank support
certificates, the finite matrix, the exact divisor criterion, and the
fixed-scale comparison kernel. A targeted literature search did not
locate these statements in this formulation. That is a limited novelty
assessment, not a proof that no equivalent result has appeared elsewhere.

\section{Hahn workspaces and support-controlled operations}
\label{sec:support}

\subsection{The scalar field and the ambient surreal interpretation}
Throughout, $\Gamma$ is a set-sized ordered abelian subgroup of
$(\No,+)$ and
\begin{equation}
 \Kg=\left\{\sum_{\gamma\in S}c_\gamma t^\gamma:
          S\subseteq\Gamma\text{ well ordered},\ c_\gamma\in\C\right\}.
 \label{eq:hahnfield}
\end{equation}
Zero coefficients are omitted from the support. We use the convention
$t^\gamma=\omega^{-\gamma}$, so positive exponents represent
infinitesimals. Surreal normal forms give an embedding
$\Kg\hookrightarrow\No[i]$; multiplication is the Hahn convolution.
For a nonzero series $f$, set $v(f)=\min\supp f$, and put
\[
 \rg=\{f:v(f)\ge0\}\cup\{0\},\qquad
 \mg=\{f:v(f)>0\}\cup\{0\}.
\]
We occasionally assume $\Gamma$ divisible to describe polynomial roots
inside $\Kg$. The root-free statements do not require divisibility.
All supports and all cohomology groups in this article are sets.

For a complex vector space $V$, the notation $V((t^\Gamma))$ means
well-ordered Hahn series with coefficients in $V$. When $V$ is
infinite-dimensional this is generally \emph{not} $V\ot\Kg$.
When $V$ is finite-dimensional it is canonically $V\ot\Kg$.
There are analogous nonnegative and positive-support versions.

\begin{definition}[Strong summability]
A family of Hahn series is strongly summable if the union of its supports
is well ordered and every exponent occurs in only finitely many members
of the family. Its sum is defined coefficientwise. No topology is part
of this definition.
\end{definition}

\subsection{The support lemma used in every infinite construction}
\begin{lemma}[Neumann support lemma and its operator form]
\label{lem:neumann}
Let $E\subset\Gamma_{>0}$ be well ordered, and let
\[
 E^*=\{e_1+\cdots+e_n:n\ge0,\ e_j\in E\}.
\]
Then $E^*$ is well ordered, and a fixed exponent has only finitely many
representations as a finite word in $E$. If $S\subset\Gamma$ is well
ordered, then $S+E^*$ is well ordered, and a fixed exponent has only
finitely many representations $s+e_1+\cdots+e_n$.

Consequently, a complex-linear coefficient operator preserves support.
If $A=\sum_{e\in E}t^e A_e$ is a positive-support operator on Hahn vectors,
then
\begin{equation}
 (1+A)^{-1}=\sum_{n\ge0}(-A)^n
 \label{eq:neumann}
\end{equation}
is a well-defined two-sided inverse. Acting on support $S$, it produces
support contained in $S+E^*$.
\end{lemma}
\begin{proof}
The assertions about positive semigroups and finite decompositions are
the classical Neumann lemma~\cite{neumann}. The assertion involving $S$
also uses the Hahn sum lemma: the sum of two well-ordered subsets of an
ordered abelian group is well ordered, with only finitely many
representations of a fixed element as a sum of one element from each.
Applying this to $S$ and $E^*$ and then to the finite word
representations proves the stated version.

For the operator assertion, a coefficient of $A^n x$ at $\gamma$ is a
sum of terms
\[
 A_{e_1}\cdots A_{e_n}(x_s),\qquad
 s+e_1+\cdots+e_n=\gamma.
\]
Only finitely many such terms occur, even after summing over $n$.
Thus every coefficient in \eqref{eq:neumann} is defined by finite
algebra. Multiplication by $1+A$ cancels the terms pairwise at each
exponent, proving both inverse identities. The displayed support bound
follows from the same representation of exponents.
\end{proof}

\begin{remark}[What the lemma does not say]
\label{rem:nonarch}
It need not be true that $n\delta$ eventually exceeds every element of
$\Gamma_{>0}$. For example, $n<\omega$ for every integer $n$ in
$\Gamma=\Q+\Q\omega$. Neither this lemma nor any later proof uses that
false cofinality assertion. Finite decomposability of an \emph{exact}
exponent, not eventual domination of every exponent, is the mechanism.
\end{remark}

\begin{corollary}[Positive exponential and logarithm]
\label{cor:explog}
In any commutative Hahn coefficient algebra, positive-support elements
admit the operations
\[
 \Exp(h)=\sum_{n\ge0}\frac{h^n}{n!},\qquad
 \Log(1+h)=\sum_{n\ge1}\frac{(-1)^{n+1}}n h^n.
\]
These are inverse group isomorphisms between the positive additive group
and the multiplicative group $1+$ positive support. If $\supp h\subset E$,
the output support is contained in $E^*$.
\end{corollary}
\begin{proof}
Lemma~\ref{lem:neumann} makes every substitution strongly summable.
The usual formal power-series identities can therefore be checked
coefficientwise, where they reduce to finite identities over $\Q$.
\end{proof}

\subsection{The exact meaning of a support certificate}
Later expressions such as $(1+h\alpha)^{-1}$ will act on spaces of smooth
forms. The coefficient operators are ordinary linear maps, such as a
Green operator or multiplication by a fixed smooth form. They need not
be uniformly bounded over all exponents; they merely have to send each
ordinary coefficient to a legitimate ordinary coefficient.
A certificate of the form
\begin{equation}
 \supp(\text{output})\subset S+E^*
 \label{eq:certificate}
\end{equation}
proves that the entire result is one Hahn series. It also explains
precisely why no exchange of a divergent analytic sum and an integral is
being made. Ordinary differentiation, integration, and Green operators
are applied to the finitely many ordinary terms contributing at each
fixed Hahn exponent.

\section{Common-domain sheaves and compact Dolbeault comparison}
\label{sec:dolbeault}

\subsection{A sheaf definition that handles disconnected opens}
\begin{definition}[The holomorphic and meromorphic Hahn sheaves]
On an ordinary open set $U\subseteq X$, a section of $\Hg$ is a family
of ordinary holomorphic coefficients $(f_\gamma)_{\gamma\in\Gamma}$
whose nonzero indices are locally contained in one well-ordered set.
The subsheaves $\Rg$ and $\Hp$ require respectively nonnegative and
positive support. The sheaf $\Mg$ is defined in the same way using
ordinary meromorphic coefficients. All coefficients of a local section
are defined on the \emph{same} ordinary open neighborhood.
\end{definition}

The local condition is essential to make this a sheaf on arbitrary
opens. On a connected nonempty $U$, the identity theorem implies the
simpler formulas
\begin{equation}
 \Hg(U)=\OO_X(U)((t^\Gamma)),\qquad
 \Mg(U)=\MM_X(U)((t^\Gamma)).
 \label{eq:connectedsections}
\end{equation}
Indeed, a nonzero holomorphic or meromorphic coefficient on connected
$U$ cannot vanish identically on a smaller nonempty open set. Hence the
support of the whole family is already bounded by the support on any
one neighborhood admitting a local bound. On a disconnected open,
different components may have different supports.

We write $\Acal^{0,q}_\Gamma$ for the analogous sheaf of smooth Hahn
$(0,q)$-forms, and use subscripts $\ge0$ and $+$ for its support
restrictions. Smooth coefficients do not satisfy the identity theorem.
However, compactness gives the uniformity needed for global analysis.

\begin{lemma}[Compact support uniformity]
\label{lem:compact}
A globally defined smooth Hahn form on compact $X$ has one well-ordered
support. The same holds for forms in an ordinary finite-rank smooth
bundle. All coefficientwise ordinary linear operators on these global
spaces are therefore well defined and support preserving.
\end{lemma}
\begin{proof}
Choose a finite subcover of neighborhoods with support bounds
$S_1,\ldots,S_N$. Every nonzero coefficient is nonzero on at least one
member, so its exponent belongs to their union. A finite union of
well-ordered subsets of a linear order is well ordered: a nonempty
subset has a least element in each nonempty intersection with $S_j$,
and the least of those finitely many elements is its least element.
For a bundle, use a finite ordinary trivializing cover; ordinary
transition matrices do not introduce exponents.
\end{proof}

\subsection{The coefficientwise resolution}
\begin{proposition}[Dolbeault resolution]
\label{prop:dolbeault}
There are fine resolutions
\[
 0\longrightarrow\Hg\longrightarrow\Acal^{0,0}_\Gamma
       \xrightarrow{\db}\Acal^{0,1}_\Gamma\longrightarrow0,
\]
and the same statement holds for the nonnegative and positive-support
versions. For an ordinary holomorphic vector bundle $V$, the analogous
resolution computes the cohomology of $V$ with Hahn coefficients.
\end{proposition}
\begin{proof}
The kernel statement is coefficientwise holomorphicity. For local
surjectivity, take a coordinate disk on which the coefficients of the
input are smooth and a fixed smaller disk. Choose one smooth cutoff
supported in the first disk and equal to one on the smaller disk.
The ordinary Cauchy--Green solution operator, with this fixed cutoff,
solves $\db$ for every coefficient on the same smaller disk. Its
coefficientwise extension preserves the support. Thus no exponent-
dependent shrinking of domains is necessary.

An ordinary smooth partition of unity acts by multiplication in every
coefficient and gives the fineness of the smooth sheaves. The bundle
statement is local in an ordinary holomorphic frame. On a curve there
are no $(0,2)$-forms, so the displayed resolution stops in degree one.
\end{proof}

Fix ordinary Hermitian metrics. Write $d=\db_V$, and let
$i_q$ include the harmonic representatives, $p_q$ be harmonic projection,
and $h$ be the ordinary Hodge homotopy from degree one to degree zero.
The needed identities are
\begin{equation}
 h d=1-i_0p_0,\qquad d h=1-i_1p_1,\qquad
 p_0h=0,\quad hi_1=0,\quad p_qi_q=1.
 \label{eq:hodge}
\end{equation}
These are ordinary compact-curve Hodge theory~\cite{narasimhan}.
By Lemma~\ref{lem:compact}, they remain valid after coefficientwise
Hahn extension, with any of our three support restrictions.

\begin{theorem}[Compact cohomology comparison]
\label{thm:comparison}
For an ordinary holomorphic vector bundle $V$ on $X$ and $q=0,1$,
\begin{align}
 H^q(X,V_\Gamma)&\simeq H^q(X,V)\ot\Kg,\label{eq:compareK}\\
 H^q(X,V_{\Gamma,\ge0})&\simeq H^q(X,V)\ot\rg,\label{eq:compareR}\\
 H^q(X,V_{\Gamma,+})&\simeq H^q(X,V)\ot\mg.\label{eq:comparem}
\end{align}
All higher cohomology groups vanish. In particular,
\[
 H^1(X,\Hp)\simeq H^1(X,\OO_X)\ot\mg,
 \qquad H^0(X,\Hg)=\Kg.
\]
\end{theorem}
\begin{proof}
The resolution identifies sheaf cohomology with the cohomology of its
global two-term complex. Extend \eqref{eq:hodge} coefficientwise.
The maps $i$ and $p$ then identify this cohomology with Hahn families
of ordinary harmonic representatives, preserving the support restriction.
Ordinary harmonic spaces are finite-dimensional, so these families are
precisely the displayed tensor products. This also shows that the
identification with ordinary cohomology is independent of the auxiliary
metrics. The last statement uses that a holomorphic function on a compact
connected Riemann surface is constant.
\end{proof}

\begin{remark}[Why the noncompact case is different]
This proof needs both finite-dimensional ordinary harmonic spaces and a
single global support bound for smooth forms. Local bounds on infinitely
many disjoint regions need not combine into a well-ordered set. The
compact result must not be applied to the noncompact plane merely by
replacing $X$ by $\C$ and setting $g=0$.
\end{remark}

\section{The complete Picard classification}
\label{sec:pic}

\subsection{Units have three independent components}
\begin{lemma}[Unit decomposition]
\label{lem:units}
There is a canonical isomorphism of abelian sheaves
\begin{equation}
 \Hg^\times\simeq\underline\Gamma\times\OO_X^\times\times\Hp,
 \qquad (\gamma,a,h)\longmapsto t^\gamma a\Exp(h).
 \label{eq:units}
\end{equation}
For the integral sheaf it restricts to
\begin{equation}
 \Rg^\times\simeq\OO_X^\times\times\Hp.
 \label{eq:integralunits}
\end{equation}
The splitting is relative to the fixed monomial section $\gamma\mapsto
 t^\gamma$ and the fixed inclusion of ordinary coefficients.
\end{lemma}
\begin{proof}
Work on a connected open set where an element $f$ and its inverse are
represented by Hahn series. Let their lowest exponents be $\gamma$ and
$\beta$, and their leading coefficients $a$ and $b$. The lowest
coefficient of the product is $ab$. Since ordinary holomorphic functions
on a connected open set form an integral domain, $\gamma+\beta=0$
and $ab=1$. Thus $a$ is a nowhere-zero ordinary holomorphic function.
There is a unique positive-support $u$ with
$f=t^\gamma a(1+u)$. Set $h=\Log(1+u)$ and apply
Corollary~\ref{cor:explog}.

Conversely, the displayed factors are invertible and their inverses
have the required supports. Uniqueness follows from the least exponent,
then the leading coefficient, then the logarithm. Multiplication adds
least exponents, multiplies leading coefficients, and adds positive
logarithms, so this is a group isomorphism. Restriction preserves the
three components by the identity theorem, making it a sheaf isomorphism.
An integral unit has both it and its inverse nonnegative, forcing
$\gamma=0$.
\end{proof}

\begin{definition}[The three classes of a line bundle]
A Hahn line bundle means a locally free rank-one $\Hg$-module. Its class
in $\Piclf(X,\Hg)=H^1(X,\Hg^\times)$ has three components under
\eqref{eq:units}: its \emph{valuation class} $\vclass(L)$, its ordinary
line-bundle component $L_0$, and its positive logarithmic class $\xi(L)$.
We define its degree to be $\deg L_0$. The ordinary component is a
projection of unit cocycles; it is not a residue map on all of $\Hg$.
\end{definition}

\begin{theorem}[Exact compact Picard classification]
\label{thm:pic}
There are canonical group isomorphisms
\begin{align}
 \Piclf(X,\Hg)
 &\simeq H^1(X,\underline\Gamma)\oplus\Pic(X)
       \oplus\bigl(H^1(X,\OO_X)\ot\mg\bigr),\label{eq:picmain}\\
 \Piclf(X,\Rg)
 &\simeq\Pic(X)\oplus\bigl(H^1(X,\OO_X)\ot\mg\bigr).
 \label{eq:picintegral}
\end{align}
Moreover,
\[
 H^1(X,\underline\Gamma)\simeq\Hom(H_1(X,\Z),\Gamma).
\]
After choosing a homology basis and a basis of $H^1(X,\OO_X)$,
\begin{equation}
 \Piclf(X,\Hg)\simeq\Gamma^{2g}\oplus\Pic(X)\oplus\mg^{\,g}.
 \label{eq:piccoordinates}
\end{equation}
Extension from $\Rg$ to $\Hg$ is injective on these Picard groups and
has image exactly the zero-valuation-class subgroup.
\end{theorem}
\begin{proof}
For a ringed space, locally free rank-one modules are classified by
unit-valued transition cocycles, modulo changes of local frames.
Taking $H^1$ of the finite direct-product decomposition in
Lemma~\ref{lem:units} gives \eqref{eq:picmain} and
\eqref{eq:picintegral}; Theorem~\ref{thm:comparison} identifies the last
summand. The ordinary unit term is the ordinary holomorphic Picard group.

A compact surface has a finite good cover. Constant-sheaf cohomology
there agrees with singular cohomology with coefficients in $\Gamma$.
The universal coefficient sequence in degree one has zero extension
term because $H_0(X,\Z)=\Z$ is free. This gives the displayed Hom group.
The first homology of a genus-$g$ surface is $\Z^{2g}$, and
$\dim_\C H^1(X,\OO_X)=g$, which proves \eqref{eq:piccoordinates}.
Finally, the map induced by integral extension includes the last two
factors and sets the valuation factor to zero.
\end{proof}

\subsection{A normal form with a finite support bound}
If $\vclass(L)=0$, take a finite good trivializing cover and change local
frames by monomials to remove the valuation cocycle. The remaining
transition functions are uniquely
\begin{equation}
 g_{ij}=a_{ij}\Exp(h_{ij}),\qquad
 a_{ij}\in\OO^\times(U_i\cap U_j),\quad h_{ij}\in\Hp(U_i\cap U_j).
 \label{eq:normalcocycle}
\end{equation}
The $a_{ij}$ define $L_0$, and the $h_{ij}$ form an additive cocycle.
A finite union of their supports gives one well-ordered
$E\subset\Gamma_{>0}$. Local smooth gauges and all later nonlinear
operations can therefore be certified in $E^*$.

\begin{remark}[The classification does not assume local stalks]
An ordinary-base Hahn stalk may not be a local ring. No statement about
all tensor-invertible sheaves on arbitrary ringed spaces is needed:
$\Piclf$ is defined by actual rank-one local trivializations. This
avoids using a local-ring hypothesis where none was established.
\end{remark}

\section{A finite matrix for the zero-valuation sector}
\label{sec:matrix}

\subsection{From holomorphic transition functions to a global operator}
\begin{proposition}[Positive Dolbeault normal form]
\label{prop:normalform}
Let $L$ have valuation class zero and transitions
\eqref{eq:normalcocycle}. There is a positive-support global scalar
$(0,1)$-form $\alpha$ such that the smooth Hahn realization of $L$ is
identified with that of $L_0$ and its holomorphic sections are the kernel
of
\begin{equation}
 D=d+\delta,\qquad d=\db_{L_0},\qquad
 \delta(s)=\alpha s.
 \label{eq:D}
\end{equation}
One may choose every coefficient of $\alpha$ harmonic. The class of
$\alpha$ is $\xi(L)$. If the supports of $h_{ij}$ lie in $E$, the
coefficients of $\alpha$ have support contained in $E$ and the smooth
gauges have support contained in $E^*$.
\end{proposition}
\begin{proof}
Fix an ordinary partition of unity $\rho_k$ subordinate to the finite
cover, and use the convention that section coordinates satisfy
$s_i=g_{ij}s_j$. Set
\[
 b_i=\sum_k\rho_k h_{ik}.
\]
Terms are extended smoothly by zero where the corresponding partition
function vanishes. The cocycle identity gives $b_i-b_j=h_{ij}$.
Thus $r_i=\Exp(-b_i)s_i$ has ordinary transition functions $a_{ij}$.
The holomorphicity equation becomes
\[
 0=\db s_i=\Exp(b_i)\bigl(\db r_i+(\db b_i)r_i\bigr).
\]
Since $\db(b_i-b_j)=0$, the forms $\db b_i$ glue to a global $\alpha$.
This is the coefficientwise Cech--Dolbeault representative of $[h_{ij}]$.
All operations so far preserve $E$, apart from exponential gauges,
which have support in $E^*$.

Every scalar $(0,1)$-form on a curve is $\db$-closed. Ordinary Hodge
splitting, applied coefficientwise, writes
$\alpha=\alpha_{\mathrm{harm}}+\db c$ with $\supp c\subset E$.
The further smooth change $r_i=\Exp(-c)r'_i$ replaces $\alpha$ by its
harmonic part. This proves all assertions.
\end{proof}

The same construction applies to an integral Hahn line bundle, using
nonnegative-support section coordinates. Local smooth solutions of
$\db c=\alpha$ and the gauge $\Exp(c)$ show that the complex
\eqref{eq:D} is its fine Dolbeault resolution. Its global smooth terms
are the Hahn extensions of the ordinary smooth terms of $L_0$.
This follows by taking a finite trivializing cover and using the support
bound for the gauges; it is not an assumption that tensor product
commutes with an infinite-dimensional space of sections.

\subsection{The ordinary Hodge data}
Let
\[
 V_0=H^0(X,L_0),\qquad V_1=H^1(X,L_0),\qquad n_q=\dim_\C V_q,
\]
and choose the Hodge maps of \eqref{eq:hodge} for $L_0$.
All maps are extended coefficientwise. On degree zero define
\begin{equation}
 T=(1+h\delta)^{-1}=\sum_{n\ge0}(-h\delta)^n,
 \qquad
 B=p_1\delta T i_0:V_0\ot\Kg\longrightarrow V_1\ot\Kg.
 \label{eq:B}
\end{equation}
Here $h\delta$ has positive support, so Lemma~\ref{lem:neumann}
applies. A useful expanded formula is
\begin{equation}
 B=\sum_{n\ge0}(-1)^n p_1\delta(h\delta)^n i_0.
 \label{eq:Bexpanded}
\end{equation}
If $\supp\alpha\subset E$, every entry of $B$ has support in
$E^*\setminus\{0\}$. In particular, the matrix is defined over $\rg$
and all its entries lie in $\mg$.

\begin{theorem}[Finite cohomology and explicit representatives]
\label{thm:finite}
For a valuation-zero Hahn line bundle $L$ represented by
\eqref{eq:D}, there are isomorphisms
\begin{equation}
 H^0(X,L)\simeq\ker B,\qquad H^1(X,L)\simeq\coker B.
 \label{eq:finitecohomology}
\end{equation}
The first is realized by $v\mapsto Ti_0v$. In degree one, a form $w$
is represented by the finite vector
\begin{equation}
 P_1w=p_1(1+\delta h)^{-1}w
      =p_1\bigl(w-\delta T h w\bigr),
 \label{eq:P1}
\end{equation}
taken modulo $\im B$.

For the normalized integral lattice $L_R$, the same statements hold
with the finite complex
\begin{equation}
 V_0\ot\rg\xrightarrow{B}V_1\ot\rg.
 \label{eq:integralcomplex}
\end{equation}
Each representative or correction constructed from input support $S$
has support contained in $S+E^*$. No valuation-rank or cofinality
hypothesis is required.
\end{theorem}
\begin{proof}
We give direct elimination in the two-term complex, so there is no
unverified appeal to an infinite-dimensional perturbation theorem.
All inverse expressions below are justified first by
Lemma~\ref{lem:neumann}. The proof uses only \eqref{eq:hodge}.

\emph{Degree zero.} Let $y=Ti_0v$. The identity
$y+h\delta y=i_0v$ gives, on applying $d$,
\[
 dy+dh\delta y=0.
\]
Since $dh=1-i_1p_1$ on degree one, this says
\begin{equation}
 Dy=i_1p_1\delta y=i_1Bv.
 \label{eq:chainI}
\end{equation}
Also $p_0y=v$, because $p_0h=0$ implies $p_0T=p_0$.
Thus a vector in $\ker B$ gives a unique section.
Conversely, if $Ds=0$, applying $h$ gives
\[
 (1+h\delta)s=i_0p_0s,
\]
so $s=Ti_0p_0s$ and \eqref{eq:chainI} implies $Bp_0s=0$.
This proves the degree-zero statement.

\emph{Reduction in degree one.} For a degree-one form $w$, put
$y=T h w$. Since $p_0T=p_0$ and $p_0h=0$, we have $p_0y=0$.
The defining equation $y+h\delta y=hw$ yields
\[
 h(w-Dy)=hw-(1-i_0p_0)y-h\delta y=0.
\]
By $dh=1-i_1p_1$, a degree-one form in the kernel of $h$ is harmonic.
It follows that
\begin{equation}
 w-DThw=i_1p_1(w-\delta Thw)=i_1P_1w.
 \label{eq:eliminate}
\end{equation}
Every class therefore has a harmonic representative. The identity
$1-\delta T h=(1+\delta h)^{-1}$ proves the other formula for $P_1$.

\emph{Relations among harmonic representatives.} Suppose $i_1b=Ds$.
Applying $h$ and using $hi_1=0$ gives
$(1+h\delta)s=i_0p_0s$, so $s=Ti_0v$ for $v=p_0s$.
Then \eqref{eq:chainI} gives $b=Bv$. Conversely, every $Bv$ is a
boundary by that same identity. Thus degree-one cohomology is the
cokernel of $B$.

Every map used here is coefficientwise linear or has the positive
Neumann support bound. The proof consequently works unchanged with
nonnegative supports over $\rg$ and gives \eqref{eq:integralcomplex}.
The same observation proves the support certificate.
\end{proof}

\begin{remark}[Relation to homological perturbation]
The operators $Ti_0$, $P_1$, and $Th$ are the usual transferred
inclusion, projection, and homotopy for a two-term contraction.
Homological perturbation and its deformation applications are
classical~\cite{crainic,greenlazarsfeld}. The point requiring proof here
is that their infinite expressions are admissible in an arbitrary-rank
Hahn workspace, and that the resulting global complex computes the
common-domain sheaf cohomology.
\end{remark}

\subsection{Riemann--Roch and stable section counts}
\begin{corollary}[Riemann--Roch in the zero-valuation sector]
\label{cor:RR}
Let $d=\deg L_0$ and $\vclass(L)=0$. Both cohomology groups are
finite-dimensional over $\Kg$, and
\begin{equation}
 h^0_\Gamma(L)-h^1_\Gamma(L)=d+1-g.
 \label{eq:RR}
\end{equation}
More precisely, for $r=\rank_{\Kg}B$,
\[
 h^0_\Gamma(L)=n_0-r,\qquad h^1_\Gamma(L)=n_1-r.
\]
In particular, $h^q_\Gamma(L)\le\dim_\C H^q(X,L_0)$.
If $d>2g-2$, then
\begin{equation}
 H^1(X,L)=0,\qquad h^0_\Gamma(L)=d+1-g.
 \label{eq:highdegree}
\end{equation}
The corresponding integral space of sections is free of rank $d+1-g$
over $\rg$, with a basis obtained by applying $Ti_0$ to an ordinary
basis of $H^0(X,L_0)$. If $d<0$, then $H^0(X,L)=0$.
\end{corollary}
\begin{proof}
The finite kernel and cokernel in Theorem~\ref{thm:finite} have the
stated dimensions. Their difference is $n_0-n_1$, which is
$d+1-g$ by ordinary Riemann--Roch. Ordinary Serre duality gives
$n_1=0$ for $d>2g-2$, so $B$ has zero target. The integral conclusion
then follows from the same finite complex. For $d<0$, ordinary
$n_0=0$, and the kernel is zero.
\end{proof}

\begin{proposition}[Degree-zero deformations of the trivial bundle]
\label{prop:degreezero}
Suppose $L_0=\OO_X$ and let the harmonic positive class be
$\xi\in H^1(X,\OO_X)\ot\mg$. Then $B:\Kg\to\Kg^g$ sends
$1$ exactly to the coordinate vector of $\xi$. Thus, if $\xi\ne0$,
\[
 h^0_\Gamma(L)=0,\qquad h^1_\Gamma(L)=g-1.
\]
For $\xi=0$ the dimensions are $1$ and $g$. On an elliptic curve,
if $\xi=\lambda\eta$ with $\eta\ne0$ an ordinary harmonic generator
and $0\ne\lambda\in\mg$, the integral groups are
\begin{equation}
 H^0(X,L_R)=0,\qquad H^1(X,L_R)\simeq\rg/\lambda\rg.
 \label{eq:elliptictorsion}
\end{equation}
\end{proposition}
\begin{proof}
The ordinary harmonic degree-zero space consists of constants. Since
$\alpha$ is harmonic, $h\delta(1)=h\alpha=0$. Hence $T(1)=1$ and
$B(1)=p_1\alpha=\xi$. A nonzero map from a one-dimensional vector space
has rank one. Over $\rg$, the elliptic complex is multiplication by
$\lambda$ on $\rg$; this ring is a domain, proving
\eqref{eq:elliptictorsion}.
\end{proof}

\subsection{Field enlargement does not remove an obstruction}
\begin{corollary}[Base extension]
\label{cor:baseextension}
For an ordered-group inclusion $\Gamma\subseteq\Gamma'$ using the same
monomials, the induced map on compact Picard groups is injective.
For a valuation-zero bundle,
\[
 H^q(X,L)\otimes_{\Kg}K_{\Gamma'}\simeq H^q(X,L_{\Gamma'}),
 \qquad q=0,1.
\]
\end{corollary}
\begin{proof}
In \eqref{eq:picmain}, both the coefficient inclusion
$\mg\hookrightarrow\mm_{\Gamma'}$ and the induced map
$\Hom(H_1(X,\Z),\Gamma)\to\Hom(H_1(X,\Z),\Gamma')$
are injective. For cohomology use the same ordinary Hodge data. The
matrix $B$ is unchanged; all its coefficients are merely viewed in a
larger field. Kernels and cokernels of finite matrices commute with a
field extension, which proves the formula.
\end{proof}

\section{Valuation monodromy is exactly the meromorphic obstruction}
\label{sec:meromorphic}

\subsection{The meromorphic convention}
A coherent meromorphic section of $L$ means a section of
$L\otimes_{\Hg}\Mg$. On a connected chart its coordinate is a Hahn
series of ordinary meromorphic functions. Individual coefficients may
have different pole orders; there is no imposed uniform bound on those
orders. In particular, this meromorphic sheaf has been specified
independently and is not silently identified with a different sheaf of
fractions or a non-Archimedean analytic function field.
Every nonzero section of $\Mg$ on a connected open set is invertible:
its leading coefficient is an invertible ordinary meromorphic function,
and its remaining factor is inverted by Lemma~\ref{lem:neumann}.

\begin{theorem}[Exact meromorphic-section criterion]
\label{thm:meromorphic}
A Hahn line bundle $L$ has a nonzero coherent meromorphic section if and
only if
\begin{equation}
 \vclass(L)=0\quad\text{in }H^1(X,\underline\Gamma).
 \label{eq:meromorphiccriterion}
\end{equation}
If $\vclass(L)\ne0$, then $H^0(X,L)=0$, regardless of its degree.
\end{theorem}
\begin{proof}
Suppose $s$ is a nonzero coherent meromorphic section. Choose a finite
connected good trivializing cover, with $s_i=g_{ij}s_j$. None of its
local coordinates can be identically zero: the coefficientwise identity
theorem propagates a zero across any nonempty overlap, and $X$ is
connected. Hence the local valuations $\beta_i=v(s_i)$ are defined and
constant on each chart. Valuations in the transition identity give
\[
 v(g_{ij})=\beta_i-\beta_j.
\]
The valuation cocycle is a coboundary, proving necessity.

Conversely, suppose $\vclass(L)=0$. Fix an ordinary point $p\in X$
and take $N\ge0$ such that $d+N>2g-2$ and $d+N+1-g>0$.
The twist $L(Np)$ still has valuation class zero. By
Corollary~\ref{cor:RR} it has a nonzero holomorphic section. The ordinary
line bundle $\OO_X(Np)$ has its usual nonzero meromorphic trivialization;
using it turns this section into a nonzero coherent meromorphic section
of $L$. Locally this only multiplies a Hahn series by an ordinary
meromorphic factor, so the common-domain and support conditions are
preserved. A nonzero holomorphic section is a meromorphic one, proving
the final assertion.
\end{proof}

\subsection{Divisor classes and the missing summand}
Define the sheaf of coherent meromorphic Cartier data by
\[
 \mathscr D_\Gamma=\Mg^\times/\Hg^\times.
\]
Set $\Div_\Gamma^{\mathrm{coh}}(X)=H^0(X,\mathscr D_\Gamma)$, and let
$\Prin_\Gamma^{\mathrm{coh}}(X)$ be the image of
$H^0(X,\Mg^\times)$. This convention includes the finite cluster
Cartier divisors of Section~\ref{sec:abel}, but does not assert that
all possible meromorphic data are finite sums of moving points.

\begin{corollary}[The divisor-class sequence]
\label{cor:divseq}
There is a split exact sequence
\begin{equation}
 0\longrightarrow
 \frac{\Div_\Gamma^{\mathrm{coh}}(X)}{\Prin_\Gamma^{\mathrm{coh}}(X)}
 \longrightarrow \Piclf(X,\Hg)
 \xrightarrow{\vclass}H^1(X,\underline\Gamma)
 \longrightarrow0.
 \label{eq:divseq}
\end{equation}
In particular,
\begin{equation}
 \frac{\Div_\Gamma^{\mathrm{coh}}(X)}{\Prin_\Gamma^{\mathrm{coh}}(X)}
 \simeq\Pic(X)\oplus\bigl(H^1(X,\OO_X)\ot\mg\bigr).
 \label{eq:divclass}
\end{equation}
\end{corollary}
\begin{proof}
The exact sequence of multiplicative sheaves defining
$\mathscr D_\Gamma$ gives an injective map from Cartier data modulo
principal data into $H^1(X,\Hg^\times)$. Its image consists exactly of
line bundles that become trivial over $\Mg$, equivalently those with a
nonzero coherent meromorphic section. By
Theorem~\ref{thm:meromorphic}, that image is $\ker\vclass$.
The monomial map $\gamma\mapsto t^\gamma$ splits the valuation map;
Theorem~\ref{thm:pic} identifies its kernel.
\end{proof}

\begin{example}[High degree without sections]
\label{ex:monodromy}
Suppose $g\ge1$ and $\Gamma\ne0$. Choose a nonzero homomorphism
$\rho:H_1(X,\Z)\to\Gamma$. A constant-sheaf cocycle for $\rho$ gives
a Hahn line bundle with monomial transitions $t^{\rho_{ij}}$.
Tensor it with any ordinary degree-$d$ line bundle. Its valuation class
remains $\rho$, so it has no nonzero coherent meromorphic or holomorphic
section, even for arbitrarily large $d$.

This does not contradict \eqref{eq:RR}, whose hypothesis is valuation
class zero. Nor does it contradict Riemann--Roch for a proper algebraic
curve over $\Kg$: the ordinary-base sheaf category is not thereby
identified with that algebraic category. A purported equivalence
preserving degree, ordinary twists, and global sections would already
fail on these examples.
\end{example}

\section{A root-free Abel theorem for moving clusters}
\label{sec:abel}

\subsection{Finite cluster data and their ordinary reduction}
Fix a finite set $A\subset X$ of distinct ordinary points. Choose
pairwise disjoint coordinate disks at the points of $A$, with coordinate
$u_c$ vanishing at $c$. On each disk take two monic polynomials
\begin{equation}
 P_c^\pm(u)=u^{m_c^\pm}
       +\sum_{r=0}^{m_c^\pm-1}a_{c,r}^\pm u^r,
       \qquad a_{c,r}^\pm\in\mg.
 \label{eq:clusters}
\end{equation}
An empty cluster has $m=0$ and $P=1$. Let $D$ denote the coherent Cartier
datum represented by $P_c^+/P_c^-$ on these disks and by $1$ off $A$.
This is legitimate: on a punctured disk each $P_c^\pm/u^{m_c^\pm}$ is a
unit of $\Hg$. Its ordinary reduction is the ordinary divisor
\begin{equation}
 D_0=\sum_{c\in A}(m_c^+-m_c^-)c.
 \label{eq:D0}
\end{equation}
The valuation class of the associated Hahn line bundle is zero.

These are divisor data at infinitesimally displaced points, not merely
values at the ordinary centers. If $\Gamma$ is divisible, the Hahn
field is algebraically closed and the roots of each nontrivial cluster
polynomial lie in $\mg$. For the latter assertion, a root with valuation
$\beta\le0$ would make its highest-power term have strictly smaller
valuation than every lower term, which cannot sum to zero. One may
therefore interpret $D$ as a finite signed multiset of infinitesimal
points in the coordinate disks. The theorem below needs no roots and
remains meaningful without this divisibility assumption.

\subsection{Newton sums are the correct symmetric coordinates}
For a monic degree-$m$ polynomial
\[
 P(u)=u^m+c_1u^{m-1}+\cdots+c_m,
\]
define its Newton power sums $p_k(P)$, for $k\ge1$, by the recurrences
\begin{align}
 p_k+c_1p_{k-1}+\cdots+c_{k-1}p_1+kc_k&=0,
                    &&1\le k\le m,\label{eq:newtonlow}\\
 p_k+c_1p_{k-1}+\cdots+c_mp_{k-m}&=0,
                    &&k>m.\label{eq:newtonhigh}
\end{align}
For $P=1$ set all $p_k=0$. These are universal symmetric polynomials
with integer coefficients; if roots are chosen, then
$p_k(P)=\sum_j\eps_j^k$, counting multiplicity.

\begin{lemma}[The logarithmic principal-part identity]
\label{lem:lognewton}
Let $E$ be the union of the supports of all nonleading coefficients in
\eqref{eq:clusters}. It is a well-ordered subset of $\Gamma_{>0}$.
For each cluster polynomial, on a punctured ordinary disk,
\begin{equation}
 \Log\frac{P(u)}{u^m}
   =-\sum_{k\ge1}\frac{p_k(P)}{k u^k}.
 \label{eq:lognewton}
\end{equation}
The right side is strongly summable as a family of
meromorphic-coefficient Hahn series, and its support is contained in
$E^*\setminus\{0\}$. At a fixed Hahn exponent, only finitely many
negative powers of $u$ occur.
\end{lemma}
\begin{proof}
The set $E$ is a finite union of well-ordered positive supports.
Writing
\[
 Q(u)=\frac{P(u)}{u^m}-1=\sum_{j=1}^m c_j u^{-j},
\]
the left side is the strongly summable series
$\sum_{n\ge1}(-1)^{n+1}Q(u)^n/n$. A fixed exponent receives contributions
from only finitely many finite words in $E$ by
Lemma~\ref{lem:neumann}. Every corresponding product is an ordinary
polynomial in $u^{-1}$ of finite degree. Therefore its coefficient at
that Hahn exponent has a finite principal part.

For independent formal coefficients, the identity
\[
 \log(1+c_1q+\cdots+c_mq^m)
          =-\sum_{k\ge1}\frac{p_k(P)}k q^k
\]
is the classical generating identity for Newton sums, obtained by
logarithmically differentiating the left side and comparing
coefficients using \eqref{eq:newtonlow}--\eqref{eq:newtonhigh}.
Substitute the positive-support $c_j$ and $q=u^{-1}$. The finite-word
argument just given licenses this substitution and the regrouping by
powers of $u^{-1}$. It proves both the identity and the stronger
summability assertion.
\end{proof}

Let $\eta$ be an ordinary holomorphic differential on $X$. At $c$ write
\begin{equation}
 \eta=\left(\sum_{n\ge0}b_{c,n}(\eta)u_c^n\right)du_c.
 \label{eq:differential}
\end{equation}
The coefficients refer to the ordinary Taylor expansion on that disk.

\begin{definition}[The infinitesimal Abel functional]
For the signed cluster datum $D$, define
\begin{equation}
 \mathcal A_D(\eta)=
 \sum_{c\in A}\sum_{k\ge1}
 \frac{b_{c,k-1}(\eta)}k
       \bigl(p_k(P_c^+)-p_k(P_c^-)\bigr).
 \label{eq:Abel}
\end{equation}
This is an element of $\mg$, linear in
$\eta\in H^0(X,\Omega_X^1)$.
\end{definition}

\begin{proposition}[Well-definedness and a residue formula]
\label{prop:Abelresidue}
The sum \eqref{eq:Abel} is strongly summable, with support contained in
$E^*\setminus\{0\}$. Put
\[
 \ell_c=\Log\frac{P_c^+}{u_c^{m_c^+}}
                  -\Log\frac{P_c^-}{u_c^{m_c^-}}.
\]
Then
\begin{equation}
 \mathcal A_D(\eta)=-\sum_{c\in A}\Res_c(\ell_c\eta).
 \label{eq:Abelresidue}
\end{equation}
For split clusters it is also the signed sum of the local primitive
displacements
\begin{equation}
 \sum_{c,j}\int_0^{\eps_{c,j}^+}\eta
       -\sum_{c,j}\int_0^{\eps_{c,j}^-}\eta,
 \qquad
 \int_0^\eps\eta:=
       \sum_{k\ge1}\frac{b_{c,k-1}(\eta)}k\eps^k.
 \label{eq:localprimitive}
\end{equation}
The integrals in this formula are Taylor--Hahn substitutions, not
fine-topological path integrals.
\end{proposition}
\begin{proof}
At a fixed Hahn exponent, Lemma~\ref{lem:lognewton} gives only finitely
many negative powers of $u_c$. Multiplying by the ordinary holomorphic
Taylor series and extracting its residue therefore uses only finitely
many Taylor coefficients. Formula~\eqref{eq:lognewton} proves
\eqref{eq:Abelresidue} and the asserted strong summability.
If roots are chosen, substitute $p_k=\sum_j\eps_j^k$.
Each primitive series is strongly summable by the positive-support
lemma, and there are finitely many roots. This proves
\eqref{eq:localprimitive}. The symmetric formula supplies the sharper
support bound in $E^*$ even when individual root supports require
fractional exponents.
\end{proof}

\subsection{The exact criterion and its proof}
\begin{theorem}[Support-controlled Abel criterion]
\label{thm:abel}
For the finite signed cluster datum $D$ of \eqref{eq:clusters}, the
following are equivalent:
\begin{enumerate}[label=\textup{(\roman*)},leftmargin=2.7em]
\item There is $F\in\Mg(X)^\times$ such that, on every cluster disk,
$F/(P_c^+/P_c^-)$ is a unit of $\Hg$, and $F$ is a unit of $\Hg$
away from $A$.
\item The ordinary divisor $D_0$ is principal and
$\mathcal A_D(\eta)=0$ for every ordinary holomorphic differential $\eta$.
\end{enumerate}
It is enough to test a basis of the $g$-dimensional space
$H^0(X,\Omega_X^1)$.

When these conditions hold, fix any ordinary meromorphic $f_0$ with
$\Div(f_0)=D_0$. One can choose $F$ with exponent-zero coefficient $f_0$
and with support contained in $E^*$. Such an $F$ is unique up to a
nonzero constant in $\Kg$. The criterion and any nonzero obstruction
are preserved under ordered-group enlargement.
\end{theorem}
\begin{proof}
Use the finite cover consisting of the cluster disks and
$U_0=X\setminus A$. The normalized line bundle
$\OO(D)\otimes\OO(-D_0)$ has positive logarithmic transition data
$\ell_c$ on the punctured disks. Its valuation class is zero and its
ordinary component is trivial. Its infinitesimal class is
\begin{equation}
 \xi_D\in H^1(X,\OO_X)\ot\mg.
 \label{eq:xiD}
\end{equation}
All coefficients of its logarithmic cocycle lie in $E^*$.

The ordinary Serre pairing between $H^1(X,\OO_X)$ and
$H^0(X,\Omega_X^1)$ can be computed by sums of residues of punctured-disk
Cech representatives. Fix the convention under which the pairing of a
Dolbeault class $[\alpha]$ with $\eta$ is
$(2\pi i)^{-1}\int_X\alpha\wedge\eta$ and the corresponding Cech
representative has value $\sum_c\Res_c(\ell_c\eta)$.
One verifies the sign by choosing a cutoff equal to one on the inner
circle and zero on the outer circle of each annulus and applying
Stokes' theorem there. This is entirely ordinary complex analysis,
applied to a single Hahn coefficient at a time.
Thus Proposition~\ref{prop:Abelresidue} gives
\begin{equation}
 \langle\xi_D,\eta\rangle=-\mathcal A_D(\eta).
 \label{eq:SerreAbel}
\end{equation}
Changing the usual Serre-pairing sign convention changes both sides'
identification by an overall sign and not its zero locus.
The ordinary pairing is perfect. Its finite-dimensional coefficientwise
extension is therefore nondegenerate on
$H^1(X,\OO_X)\ot\mg$. Consequently,
\[
 \xi_D=0\quad\Longleftrightarrow\quad
      \mathcal A_D(\eta)=0\text{ for every }\eta.
\]

If (i) holds, the line-bundle class of $D$ vanishes. Its ordinary
component is $[D_0]$ and its positive component is $\xi_D$.
The Picard classification therefore gives (ii).
Conversely, (ii) annihilates these two components, while the valuation
component is already zero. The divisor-class sequence
\eqref{eq:divseq} implies principality, which is exactly (i).

We prove the extra support and normalization assertions constructively.
When $\xi_D=0$, every coefficient of the logarithmic Cech cocycle is an
ordinary coboundary on this fixed cover. A support-preserving splitting
can be made explicitly: use a fixed partition of unity to obtain smooth
$b_i$ with $b_i-b_j=\ell_{ij}$, form the global exact $(0,1)$-form
$\alpha=\db b_i$, and apply the ordinary Hodge homotopy to set
$c=h\alpha$. Then $k_i=b_i-c$ is holomorphic and
$k_i-k_j=\ell_{ij}$. Apply this coefficientwise. No domains are shrunk,
and $\supp k_i\subset E^*\setminus\{0\}$.

Let $f_i$ be the local defining function for $D$, and let $f_i^0$ be
$u_c^{m_c^+-m_c^-}$ on a cluster disk and $1$ on $U_0$.
The functions
\[
 q_i=\frac{f_i}{f_i^0}\Exp(-k_i)
\]
agree on overlaps, have meromorphic coefficients and exponent-zero
coefficient one, and have support contained in $E^*$. They define
$q\in\Mg(X)$. Then $F=f_0q$ has the desired local unit ratios and
leading coefficient $f_0$.

Two such generators have quotient a global unit of $\Hg$.
Theorem~\ref{thm:comparison} gives $\Hg(X)=\Kg$, so that quotient is
in $\Kg^\times$. Finally, ordinary principality does not change under
workspace enlargement, and nonzero Hahn coefficients of
$\mathcal A_D$ remain nonzero. This proves the last assertion.
\end{proof}

\begin{remark}[Ordinary periods versus infinitesimal displacement]
The requirement that $D_0$ be principal is the ordinary Abel condition:
it includes degree zero and the usual condition in the complex
Jacobian~\cite{narasimhan}. The additional functional
$\mathcal A_D$ has only positive Hahn exponents. An ordinary nonzero
period lies at exponent zero, so it cannot cancel a positive-support
obstruction. There is no unmentioned period quotient in
\eqref{eq:Abel}; the infinitesimal condition is exact equality to zero.
\end{remark}

\begin{remark}[Coordinate independence]
For fixed moving-point data, replacing local coordinates changes their
polynomial descriptions but not the sum of the local primitive
displacements of $\eta$. Equivalently, the class $\xi_D$ and its Serre
pairing are intrinsic. Thus the vanishing criterion is coordinate
independent. A polynomial such as $u^2-s$ itself describes a
coordinate-dependent motion; changing its coefficients without
transporting the points is a change of divisor, not a contradiction.
\end{remark}

\section{Examples: genus zero, elliptic curves, and a nonlinear obstruction}
\label{sec:examples}

\subsection{The sphere}
If $X=\mathbb P^1(\C)$, then $g=0$ and
\[
 \Piclf(X,\Hg)\simeq\Pic(X)\simeq\Z.
\]
There is neither valuation monodromy nor an infinitesimal Abel
obstruction. Every finite signed cluster datum of total degree zero is
principal. In one affine coordinate, products of the defining
polynomials give the expected rational generator, with the factor at
infinity accounting for the total degree.

This is not a global Mittag--Leffler assertion for the noncompact
ordinary plane. There are only finitely many prescribed clusters on a
compact base, and the support union is automatically well ordered.

\subsection{An elliptic curve: exact displacement conservation}
Let $X=\C/\Lambda$ be an ordinary elliptic curve and use local
coordinates induced by its uniformizing variable. A generator of
$H^0(X,\Omega_X^1)$ is $dz$. In these coordinates
$b_{c,0}=1$ and $b_{c,n}=0$ for $n>0$, so
\begin{equation}
 \mathcal A_D(dz)=\sum_c\bigl(p_1(P_c^+)-p_1(P_c^-)\bigr).
 \label{eq:ellipticAbel}
\end{equation}
The test is exact conservation of the sum of infinitesimal
displacements, in addition to the ordinary Abel condition.

For $P^+(u)=u-\eps$ and $P^-(u)=u$ at one center, with
$0\ne\eps\in\mg$, the ordinary divisor is zero but
$\mathcal A_D(dz)=\eps$. Hence a single moved point minus its original
point is not principal. In contrast,
$P^+(u)=u^2-s$ and $P^-(u)=u^2$ have $p_1=0$ for both clusters,
so that signed divisor is principal in the uniformizing coordinate.

The same elliptic curve also gives a particularly transparent
cohomology example. An infinitesimal nontrivial degree-zero line bundle
has no cohomology over $\Kg$, while its normalized integral model has
$H^1=\rg/\lambda\rg$ as in \eqref{eq:elliptictorsion}. Passing to the
field discards that torsion; reducing to ordinary coefficients discards
the infinitesimal class. These are different operations.

\subsection{A genus-two cluster with zero first displacement but nonzero Abel class}
Consider the smooth projective completion of
\begin{equation}
 y^2=1+x^5.
 \label{eq:hyperelliptic}
\end{equation}
The polynomial $1+x^5$ has five simple roots. The double cover of the
$x$-sphere is therefore branched at those five points and at infinity,
so Riemann--Hurwitz gives genus two. A basis of holomorphic
differentials is
\begin{equation}
 \eta_0=\frac{dx}{y},\qquad\eta_1=\frac{x\,dx}{y}.
 \label{eq:genus2forms}
\end{equation}
These are holomorphic at the finite branch points in their square-root
parameters and at infinity in the parameter $x^{-1/2}$; their linear
independence and the genus calculation show that they form a basis.

At $c=(0,1)$ use $u=x$ and the branch
$y=(1+u^5)^{1/2}$. Fix $s=t^\gamma$ with $\gamma>0$ and take
\begin{equation}
 P^+(u)=u^2-s,\qquad P^-(u)=u^2.
 \label{eq:balancedcluster}
\end{equation}
The ordinary divisor is zero. In a divisible workspace the two positive
roots are $\sqrt{s}$ and $-\sqrt{s}$, so their first power sum vanishes.
Nevertheless this does not make the Abel class vanish on a genus-two
curve.

The power sums of $P^+$ are
\[
 p_{2r}=2s^r,\qquad p_{2r+1}=0,
\]
while all positive power sums of $P^-$ vanish. Since
\[
 (1+u^5)^{-1/2}=\sum_{n\ge0}\binom{-1/2}{n}u^{5n},
\]
formula \eqref{eq:Abel} gives the exact Hahn identities
\begin{align}
 \mathcal A_D(\eta_0)
 &=\sum_{\ell\ge0}
   \frac{\binom{-1/2}{2\ell+1}}{5\ell+3}s^{5\ell+3}
   =-\frac{s^3}{6}-\frac{5s^8}{128}+\cdots,\label{eq:genus2A0}\\
 \mathcal A_D(\eta_1)
 &=\sum_{\ell\ge0}
   \frac{\binom{-1/2}{2\ell}}{5\ell+1}s^{5\ell+1}
   =s+\frac{s^6}{16}+\frac{35s^{11}}{1408}+\cdots.
 \label{eq:genus2A1}
\end{align}
The leading term $s$ in \eqref{eq:genus2A1} is nonzero, so
Theorem~\ref{thm:abel} proves that the balanced cluster
\eqref{eq:balancedcluster} is \emph{not} principal.
The first symmetric displacement is zero, but a nonlinear Abel
obstruction remains. All terms belong to the monoid $\{n\gamma:n\ge0\}$;
no claim that $n\gamma$ is cofinal in $\Gamma$ is needed.

\subsection{Why the example is not a numerical approximation}
The displayed leading nonzero coefficient already certifies
nonprincipality. An arbitrary extension of the value group cannot
remove it. The higher terms describe the full obstruction, rather
than an error term in the fine topology of $\No[i]$.
The supplied program verifies the Newton identities and the first
coefficients of \eqref{eq:genus2A0}--\eqref{eq:genus2A1} exactly over
$\Q$; the infinite formulas and their implication for divisors are
proved above.

\section{The exact information lost by one-scale completion}
\label{sec:adic}

The integral sheaf is essential in this section. In $\Hg$, every
$t^\delta$ is a unit, so quotienting by $t^{n\delta}\Hg$ would give
the zero sheaf. Instead we study
\[
 \Rg/t^{n\delta}\Rg,\qquad \delta\in\Gamma_{>0}.
\]
Even its positive elements need not be nilpotent: if $0<\eps\ll\delta$,
then $t^{k\eps}$ can survive for every integer $k$. Logarithms and
exponentials are still defined by strong summability, not by falsely
assuming that every such series terminates.

\subsection{Truncated coefficients and their Picard groups}
Write
\[
 J_n=t^{n\delta}\rg,\qquad
 \RR_n=\Rg/t^{n\delta}\Rg,\qquad
 \mm_n=\mg/J_n.
\]
Here $J_n\subset\mg$. Coefficients of $\RR_n$ may be represented
uniquely by Hahn series with exponents in $[0,n\delta)$.

\begin{lemma}[Truncated Picard classification]
\label{lem:truncatedpic}
For each positive integer $n$,
\begin{equation}
 \Piclf(X,\RR_n)\simeq\Pic(X)
       \oplus\bigl(H^1(X,\OO_X)\ot\mm_n\bigr).
 \label{eq:truncatedpic}
\end{equation}
Under this identification the reduction map from
$\Piclf(X,\Rg)$ is the identity on $\Pic(X)$ and coefficient
reduction on the second summand.
\end{lemma}
\begin{proof}
Cutting off the exponents $\ge n\delta$ gives a canonical additive
representative of a local quotient section. Local representatives glue
coefficientwise; on a connected ordinary open the identity theorem
ensures that their support is well ordered. Thus the quotient sheaf
really is the indicated sheaf of truncated Hahn coefficients, not just
a pointwise quotient notation.

A truncated unit has invertible exponent-zero coefficient. To see
necessity, the constant coefficients of two inverse representatives
multiply to one. For sufficiency, divide by this ordinary coefficient
and use the geometric Hahn inverse before truncating. Positive
exponential and logarithm respect the ideal $J_n$, because adding a
nonnegative exponent to one at least $n\delta$ cannot return below
$n\delta$. They therefore descend to inverse group maps on the
positive truncated sheaf. This proves the unit decomposition
\[
 \RR_n^\times\simeq\OO_X^\times\times\PP_n,
\]
where $\PP_n$ is the positive truncated additive sheaf.

The proof of the Dolbeault comparison applies with exponents restricted
to $(0,n\delta)$. The fixed local solution operators and the global
Hodge maps preserve that restriction, and compactness again gives one
support bound. Hence
$H^1(X,\PP_n)=H^1(X,\OO_X)\ot\mm_n$.
Taking $H^1$ of the unit decomposition proves the formula and its
compatibility with reduction.
\end{proof}

\subsection{The visible convex subgroup}
Define
\begin{equation}
 \Delta=\{\gamma\in\Gamma:|\gamma|\le n\delta
                         \text{ for some integer }n\ge1\},
 \qquad I_\delta=\bigcap_{n\ge1}t^{n\delta}\rg.
 \label{eq:DeltaI}
\end{equation}
Then $\Delta$ is the convex subgroup of $\Gamma$ generated by
$\delta$. Let $R_\Delta=\C((t^\Delta))_{\ge0}$ and
$\mm_\Delta=\C((t^\Delta))_{>0}$. Truncation to exponents in $\Delta$
defines a split surjective ring map
\begin{equation}
 \rg\longrightarrow R_\Delta.
 \label{eq:convexcut}
\end{equation}
Its kernel is $I_\delta$: a nonzero element is in every $J_n$ exactly
when its least exponent is above all integer multiples of $\delta$.
Convexity ensures that all its other exponents then also lie outside
$\Delta$. Conversely, all positive exponents outside $\Delta$ exceed
every $n\delta$.

\begin{lemma}[The inverse limit of the coefficient truncations]
\label{lem:coefficientlimit}
There are canonical isomorphisms
\begin{equation}
 \varprojlim_n\rg/J_n\simeq R_\Delta,
 \qquad
 \varprojlim_n\mg/J_n\simeq\mm_\Delta.
 \label{eq:coefficientlimit}
\end{equation}
\end{lemma}
\begin{proof}
A compatible family determines each coefficient with exponent in
$\Delta_{\ge0}$ by taking a sufficiently large cutoff. Let $S$ be the
union of its nonzero indices. We check that $S$ is well ordered.
For any nonempty subset $T\subset S$, choose $\gamma_0\in T$ and then
$n$ with $\gamma_0<n\delta$. The nonempty subset
$T\cap[0,n\delta)$ lies in the well-ordered support of the $n$th
truncation and has a least element. That element is also least in all
of $T$, because the other exponents are at least $n\delta$.
Thus the compatible family defines an element of $R_\Delta$.

Restriction of a Hahn series in $R_\Delta$ gives the inverse operation.
The maps respect addition and multiplication, since a coefficient
below $n\delta$ in a product of nonnegative series only uses exponents
below that cutoff. Omitting the exponent-zero coefficient proves the
positive version.
\end{proof}

\begin{theorem}[Exact one-scale Picard comparison]
\label{thm:adic}
Let $\delta>0$ and define $\Delta,I_\delta$ by
\eqref{eq:DeltaI}. The natural map
\begin{equation}
 \Piclf(X,\Rg)\longrightarrow
       \varprojlim_n\Piclf(X,\Rg/t^{n\delta}\Rg)
 \label{eq:piccompletion}
\end{equation}
is split surjective. Its target and kernel are respectively
\begin{align}
 \varprojlim_n\Piclf(X,\Rg/t^{n\delta}\Rg)
    &\simeq\Pic(X)\oplus
             \bigl(H^1(X,\OO_X)\ot\mm_\Delta\bigr),
        \label{eq:piclimtarget}\\
 \ker\eqref{eq:piccompletion}
    &\simeq H^1(X,\OO_X)\ot I_\delta.
        \label{eq:piclimkernel}
\end{align}
The map is injective if and only if $g=0$ or $\Delta=\Gamma$.
\end{theorem}
\begin{proof}
Apply Lemma~\ref{lem:truncatedpic} at every level. The ordinary Picard
factor has identity transition maps, and the other factor is a tensor
product with the fixed finite-dimensional vector space
$H^1(X,\OO_X)$. After choosing a finite basis, this is a finite direct
power, which commutes with the inverse limit. Lemma~\ref{lem:coefficientlimit}
therefore proves \eqref{eq:piclimtarget}.

Under the full integral Picard classification, the map is the identity
on $\Pic(X)$ and the convex-support cutoff
$\mg\to\mm_\Delta$ on each infinitesimal coordinate. Inclusion of
$\Delta$-supported coefficients provides a splitting, and its kernel
is precisely \eqref{eq:piclimkernel}.

If $g=0$ the kernel is zero. For $g>0$, it is zero exactly when
$I_\delta=0$. If $\Delta=\Gamma$, no nonzero series can have least
exponent above every $n\delta$, so $I_\delta=0$. If
$\Delta\ne\Gamma$, choose a positive exponent $\eta\notin\Delta$.
Then $t^\eta\in I_\delta$ is nonzero. This proves the equivalence.
\end{proof}

\begin{remark}[What is, and is not, being completed]
The theorem computes an inverse limit of \emph{isomorphism classes of
line bundles} in the indicated integral quotient sheaves. It does not
replace them by bundles over the full surcomplex fine topology, and it
does not infer a sheaf equivalence from a formal slogan about completion.
The proof is an explicit compatible-coefficient calculation. Nor is the
source Picard group a priori assigned any complete Hausdorff topology.
The kernel proves that a topology based on powers of a single scale
would be nonseparated on its infinitesimal part whenever
$g>0$ and $\Delta\ne\Gamma$.
\end{remark}

\subsection{A deformation beyond every integer-scale jet}
\begin{example}[The $t^\omega$ witness]
\label{ex:invisible}
Take $\Gamma=\Q+\Q\omega\subset\No$, with its inherited order,
$\delta=1$, and a compact curve with $g\ge1$. Then
$\Delta=\Q$ and $t^\omega\in I_1$.
Choose a nonzero class $\zeta\in H^1(X,\OO_X)$, represented on a finite
cover by an ordinary additive cocycle $h_{ij}$. The transitions
\begin{equation}
 g_{ij}=\Exp(t^\omega h_{ij})
 \label{eq:invisiblebundle}
\end{equation}
define an integral line bundle with class $t^\omega\zeta\ne0$.
It is nontrivial both over $\Rg$ and, by Theorem~\ref{thm:pic},
over $\Hg$. But every transition reduces to $1$ modulo $t^n\Rg$
for every positive integer $n$, so every one of those reductions is
trivial.

For an elliptic curve, the integral cohomology is
$H^1\simeq\rg/t^\omega\rg$ and the full Hahn cohomology is zero in
both degrees. The nontriviality is not a numerical failure to compute
enough terms: no integer multiple of the first scale reaches the
obstruction. A larger exponent scale is mathematically necessary.
\end{example}

\subsection{No set of cutoffs detects every surreal workspace}
\begin{corollary}[A set-indexed detection obstruction]
\label{cor:setcutoffs}
Let $X$ have $g\ge1$, and let $B$ be any set of positive surreal
exponents. There exist a set-sized ordered subgroup
$\Gamma\subset\No$ containing $B$ and a nontrivial integral Hahn line
bundle $L_R$ such that
\[
 L_R\bmod t^\beta\Rg\text{ is trivial for every }\beta\in B.
\]
The same can be required for every positive integer multiple of every
$\beta\in B$.
\end{corollary}
\begin{proof}
The set of all such integer multiples is still a set. Every set of
surreal numbers has a surreal strict upper bound, for example one
obtained by the usual set-sized cut construction. Choose a positive
$\eta$ strictly above that set, and let $\Gamma$ be the divisible
additive subgroup generated by $B\cup\{\eta\}$. This is set-sized.
Use the cocycle $\Exp(t^\eta h_{ij})$ for a nonzero ordinary class
$[h_{ij}]$. Its Picard class is nonzero, while all the stated
reductions are trivial because $\eta>n\beta$. No cohomology on a
proper-class space is formed anywhere in this construction.
\end{proof}

\section{Consequences, limitations, and further mathematical questions}
\label{sec:conclusions}

\subsection{Three tests, not one}
The results separate three otherwise easily conflated decisions.
To determine whether a Hahn line bundle has any nonzero coherent
meromorphic section, compute its valuation class. To determine whether
a finite moving divisor is principal, compute the ordinary divisor
class and its positive Abel functional. To determine whether a
one-scale formal test is faithful, determine whether that scale is
cofinal in the value group.

None of these tests subsumes the others. A nonzero valuation class
prevents all meromorphic sections, even in high degree. A nonzero
infinitesimal class can have many meromorphic sections while having no
nowhere-vanishing holomorphic section. A nonzero infinitesimal class
can also be invisible to every power of a smaller scale. The genus-two
example additionally shows that cancellation of a first displacement
does not make the full infinitesimal Abel class vanish.

\subsection{Why arbitrary-rank support control is substantive}
In a rank-one complete valued field one often proves an inverse formula
by showing that the terms approach zero rapidly. That argument is not
available for all of $\Gamma$. For example, a term supported at
$\omega$ cannot be reached by multiplying the first scale $1$ by
ordinary integers. The present proofs instead use exact-exponent
finiteness. This is why the same inverse matrix and residue formulas
are valid both at ordinary small scales and at arbitrarily remote
surreal scales, while their detection by one chosen adic system may
fail.

Compactness supplies a second, independent form of control. It permits
a finite smooth cover and a uniform well-ordered support. The ordinary
Hodge projection then has finite-dimensional image. Without the first
condition, a partition-of-unity construction can create a forbidden
union of supports; without the second, a finite matrix need not result.
Both are used explicitly in the proof, rather than folded into a vague
notion of analytic convergence.

\subsection{What this article does not identify}
The objects studied are not compact subsets of the fine surcomplex
plane. The base topology is ordinary complex topology, and the Hahn
field enters the coefficients. Nor have we identified this ringed
space with an algebraic curve over $\Kg$, a rigid analytic curve,
a Berkovich curve, or a logarithmic curve. Such identifications would
require comparison functors with stated hypotheses. The valuation
monodromy examples already rule out the most naive comparison that
would preserve both degree and all section spaces.

The Abel theorem concerns finite polynomial cluster Cartier data. It
does not assert that every section of
$\mathscr D_\Gamma$ has such a presentation, nor assign a classical
finite zero count to arbitrary meromorphic-coefficient Hahn series.
Similarly, cohomology in the nonzero-valuation sector was not computed
by the positive perturbation argument: there the initial global
ordinary bundle and scalar positive gauge do not exist in the form
required by Proposition~\ref{prop:normalform}. The proved vanishing
of sections in that sector does not, by itself, determine $H^1$ or an
Euler characteristic.

\subsection{Specific next questions suggested by the proofs}
A natural continuation is to compute $H^1$ for a pure monomial local
system $t^\rho$ with $\rho\ne0$. The ordinary-base common-support
condition, rather than the abstract local system alone, must enter
such a computation. It would clarify the precise cohomological cost
of retaining the additional $\Gamma^{2g}$ factor.

Another is a comparison of the zero-valuation sector with a suitable
algebraic or non-Archimedean analytic Picard functor for the same
ordinary curve after scalar extension. The exact split classification
makes the source explicit, but does not replace the work needed to
construct and prove full faithfulness of a comparison functor.

Finally, the root-free Abel functional suggests a certified symbolic
procedure for curves whose local differentials and ordinary divisor
classes are computable. The present formulas make every fixed Hahn
coefficient a finite computation. A uniform decision procedure for
arbitrary symbolic coefficient families would require additional
representability and equality assumptions; no such algorithm is
asserted for unrestricted surcomplex inputs.

\appendix
\section{A proof and dependency audit}
\label{app:audit}

The major assertions can be checked in the following dependency order.
The table distinguishes classical inputs from the statements established
for the present Hahn category.

\begin{center}
\small
\begin{tabular}{@{}p{0.27\linewidth}p{0.65\linewidth}@{}}
\toprule
Statement & Inputs and load-bearing hypotheses\\
\midrule
Compact comparison & Ordinary Dolbeault and Hodge theory; common domains;
compactness; coefficientwise support preservation.\\[3pt]
Picard splitting & Explicit unit decomposition; classification by
unit cocycles; compact comparison; ordinary surface homology.\\[3pt]
Finite obstruction matrix & Positive Dolbeault normal form;
Neumann finite-word lemma; the two Hodge identities; finite-dimensional
ordinary harmonic spaces.\\[3pt]
Meromorphic criterion & Coefficientwise identity theorem for necessity;
finite-matrix Riemann--Roch after an ordinary high-degree twist for
sufficiency.\\[3pt]
Moving-cluster Abel test & Newton logarithm identity with a symmetric
support certificate; ordinary Serre duality; the Picard splitting.\\[3pt]
One-scale Picard kernel & Truncated unit splitting; compact cohomology;
well-ordering of the compatible support union in the convex subgroup
$\Delta$.\\
\bottomrule
\end{tabular}
\end{center}

\paragraph{Potential failure points explicitly avoided.}
An infinite-dimensional smooth Hahn space was never replaced by an
algebraic tensor product. Compactness was used before applying a global
Green operator to a Hahn family. Meromorphic coefficient pole orders
were allowed to increase with the exponent, but the ordinary domain
was fixed. The exponential was used only for positive support. No
claim that $n\delta$ is cofinal was used unless $\Delta=\Gamma$ had
already been imposed. The ideal quotients of Section~\ref{sec:adic}
were formed in the integral sheaf, not the full field sheaf. Finally,
classical ordinary results were used only on ordinary complex curves;
no unproved general surreal version of Hodge theory or Serre duality
was assumed.

\paragraph{A sharper normalization remark.}
In Theorem~\ref{thm:abel}, the constructed coefficients of $F$ are
holomorphic away from the finite set $A$, since that is true on the
fixed open $U_0=X\setminus A$. They are therefore ordinary meromorphic
functions whose poles occur only in $A$, though their pole orders may
increase. If two generators both have exponent-zero coefficient $f_0$
and nonnegative support, their constant quotient is in $1+\mg$,
a refinement of the stated uniqueness up to $\Kg^\times$.

\paragraph{Classical facts used without a fresh foundational proof.}
The ordinary Hodge decomposition, ordinary Riemann--Roch and Serre
duality, the classification of holomorphic line bundles by
$H^1(\OO^\times)$, the homology of a compact genus-$g$ surface, and the
Neumann support lemma are standard inputs. All transfers from those
facts to the specified Hahn category are proved in the article.
Consequently the manuscript is self-contained at the level of those
explicitly named foundations, rather than a replacement for an entire
textbook on compact Riemann surfaces or generalized series.

\section{Exact symbolic verification and reproducibility}
\label{app:checks}

The accompanying \texttt{verification.py} is an ordinary Python program
using SymPy. Running it produces \texttt{verification\_results.json}.
The recorded execution passed \textbf{53 of 53 exact checks}. The suite
has no floating-point tolerances and makes no internet requests.

It checks the Newton/logarithm identity for monic polynomials of degrees
two and three through power nine, verifies the displayed initial
coefficients of the genus-two Abel functionals, and verifies the
finite-dimensional contraction identities underlying
Theorem~\ref{thm:finite}. The last test uses the following purely
algebraic two-term model, not an assertion that these matrices arise
from a particular curve:
\[
 d=h=\begin{pmatrix}0&0&0\\0&1&0\\0&0&1\end{pmatrix},\qquad
 i=\begin{pmatrix}1\\0\\0\end{pmatrix},\qquad p=i^{\mathsf T},
\]
\[
 \delta=t\begin{pmatrix}0&1&2\\3&1&0\\1&-1&2\end{pmatrix}
       +t^2\begin{pmatrix}1&0&-1\\0&2&1\\1&0&0\end{pmatrix}.
\]
Here the transferred $1\times1$ matrix is exactly
\begin{equation}
 B(t)=\frac{t^2(2t^4+5t^3+4t^2-12t-4)}
                  {5t^3+4t^2+3t+1}.
 \label{eq:testB}
\end{equation}
The program verifies equality with the Schur complement and the
identities
\[
 D I_0=iB,\qquad HD=1-I_0p,\qquad
 DH=1-iP_1,\qquad P_1D=Bp,
\]
where $I_0=(1+h\delta)^{-1}i$, $H=(1+h\delta)^{-1}h$, and
$P_1=p(1+\delta h)^{-1}$. It also compares the first five positive
coefficients with the Neumann expansion. The leading order $t^2$ in
\eqref{eq:testB} illustrates why a linearized obstruction need not be
the entire transferred matrix in a general two-term deformation.

These finite certificates do not prove the infinite support lemma,
compact sheaf comparison, Picard classification, or the nonexistence
of a meromorphic section. Those assertions rest on the proofs in the
article. The program is neither a Lean formalization nor a novelty
checker.

\subsection*{Rebuilding}
From the package directory run
\begin{verbatim}
python verification.py
latexmk -pdf -interaction=nonstopmode -halt-on-error article.tex
\end{verbatim}
The bibliography is embedded in this file; no external bibliography
processor or image assets are needed. Two or three direct
\texttt{pdflatex} runs also suffice to resolve cross-references.
The archive includes a \texttt{Makefile} and a short provenance and
scope record in \texttt{README.md}.

\begin{thebibliography}{99}
\raggedright
\bibitem{repo}
V. Reshetnikov, \emph{Surreal}, GitHub repository, documentation snapshot
\texttt{aa846271b4dcae2c055b216126a87210292ec19b}.
The documentation index and the global-divisors README were inspected.
\url{https://github.com/VladimirReshetnikov/Surreal/tree/aa846271b4dcae2c055b216126a87210292ec19b/docs}.
Research drafts; not treated as independently verified theorems here.

\bibitem{neumann}
B. H. Neumann, \emph{On ordered division rings},
Transactions of the American Mathematical Society \textbf{66} (1949),
202--252. \url{https://doi.org/10.1090/S0002-9947-1949-0032593-5}.

\bibitem{gonshor}
H. Gonshor, \emph{An Introduction to the Theory of Surreal Numbers},
London Mathematical Society Lecture Note Series \textbf{110},
Cambridge University Press, 1986.

\bibitem{mantova}
V. Mantova and M. Matusinski, \emph{Surreal numbers with derivation,
Hardy fields and transseries: a survey}, in \emph{Ordered Algebraic
Structures and Related Topics}, Contemporary Mathematics \textbf{697},
American Mathematical Society, 2017, 265--290.
\url{https://arxiv.org/abs/1608.03413}.

\bibitem{narasimhan}
R. Narasimhan, \emph{Compact Riemann Surfaces}, Lectures in Mathematics
ETH Z\"urich, Birkh\"auser, 1992.
See the chapters on Dolbeault's lemma, Serre duality,
Riemann--Roch, and the Jacobian and Abel's theorem.
\url{https://doi.org/10.1007/978-3-0348-8617-8}.

\bibitem{crainic}
M. Crainic, \emph{On the perturbation lemma, and deformations},
2004, arXiv:math/0403266.
\url{https://arxiv.org/abs/math/0403266}.

\bibitem{greenlazarsfeld}
M. Green and R. Lazarsfeld,
\emph{Higher obstructions to deforming cohomology groups of line bundles},
Journal of the American Mathematical Society \textbf{4} (1991), 87--103.
\url{https://doi.org/10.1090/S0894-0347-1991-1076513-1}.

\bibitem{muller}
J. M\"uller and A. Strohmaier,
\emph{The theory of Hahn meromorphic functions, a holomorphic Fredholm
theorem and its applications}, Analysis \& PDE \textbf{7} (2014),
745--770. \url{https://doi.org/10.2140/apde.2014.7.745}.
\url{https://arxiv.org/abs/1205.0236}.

\bibitem{molchowise}
S. Molcho and J. Wise,
\emph{The logarithmic Picard group and its tropicalization},
arXiv:1807.11364, version 5.
\url{https://arxiv.org/abs/1807.11364}.
The logarithmic-curve category and its bounded-monodromy condition
are distinct from the ordinary-base Hahn category considered here.
\end{thebibliography}
\end{document}
